% Important questions only; all explanations live in their theory topics.
\begin{multicols}{2}\begin{enumerate}
\item ডেটা কমিউনিকেশন কী?
\item ডেটা কমিউনিকেশন ব্যবস্থার পাঁচটি উপাদানের নাম লেখো।
\item যোগাযোগে প্রোটোকল কেন প্রয়োজন?
\item ব্যান্ডউইথ কী?
\item Simplex, Half-duplex ও Full-duplex মোডের পার্থক্য লেখো।
\item প্রাপকভিত্তিক ডেটা পাঠানোর পদ্ধতিগুলোর পার্থক্য লেখো।
\item নির্দেশিত ও অনির্দেশিত মাধ্যম কী?
\item টুইস্টেড পেয়ার কেবলের বৈশিষ্ট্য লেখো।
\item অপটিক্যাল ফাইবার গুরুত্বপূর্ণ কেন?
\item রেডিও তরঙ্গ ও মাইক্রোওয়েভের বৈশিষ্ট্য তুলনা করো।
\item ইনফ্রারেড কোথায় ব্যবহৃত হয়?
\item স্যাটেলাইট যোগাযোগ কী?
\item Bluetooth, Wi-Fi ও WiMAX-এর ব্যবহার তুলনা করো।
\item হ্যান্ডঅফ কী?
\item রোমিং কী?
\item কম্পিউটার নেটওয়ার্ক কী?
\item LAN, MAN ও WAN-এর পার্থক্য লেখো।
\item Client–server নেটওয়ার্ক কী?
\item Peer-to-peer নেটওয়ার্ক কী?
\item Internet, Intranet ও Extranet তুলনা করো।
\item NIC-এর কাজ কী?
\item মডেমের কাজ কী?
\item হাব ও সুইচের কাজ তুলনা করো।
\item রাউটার কীভাবে নেটওয়ার্ক সংযুক্ত করে?
\item রিপিটার ও ব্রিজের কাজ লেখো।
\item গেটওয়ে কেন প্রয়োজন?
\item ফায়ারওয়াল কেন গুরুত্বপূর্ণ?
\item নেটওয়ার্ক টপোলজি কী?
\item Bus ও Star বিন্যাস তুলনা করো।
\item Mesh-এ নির্ভরযোগ্যতা বাড়তে পারে কেন?
\item Tree কোথায় ব্যবহারযোগ্য?
\item Hybrid Topology কী?
\item ক্লাউড কম্পিউটিং কী?
\item ক্লাউড স্টোরেজ কী?
\item IaaS, PaaS ও SaaS-এর পার্থক্য লেখো।
\item Public ও Private Cloud তুলনা করো।
\item Hybrid Cloud কেন ব্যবহার করা হয়?
\item ক্লাউডের সুবিধা লেখো।
\item ক্লাউডের সীমাবদ্ধতা লেখো।
\item নিরাপদে ক্লাউড ব্যবহারের নিয়মগুলো লেখো।
\item Serial ও Parallel Transmission-এর পার্থক্য কী?
\item বিট সিনক্রোনাইজেশন কেন প্রয়োজন?
\item Asynchronous ক্যারেক্টারের স্টার্ট ও স্টপ বিটের কাজ কী?
\item বাফারে ডেটা জমিয়ে Synchronous ব্লক পাঠানোর সুবিধা ও সীমা কী?
\item Isochronous সময়ের প্রয়োজন কোথায় দেখা যায়?
\item স্টার্ট, ডেটা, প্যারিটি ও স্টপ বিট দেওয়া থাকলে মোট বিট ও কার্যদক্ষতা কীভাবে নির্ণয় করা যায়?
\item পথ, সময়ের সমন্বয়, দিক ও প্রাপকভিত্তিক শ্রেণি পরস্পরের বিকল্প নয় কেন?
\item বেশি ব্যান্ডউইথ থাকলেই শব্দ নিরবচ্ছিন্নভাবে পৌঁছায় কি? Throughput, Latency ও Jitter-এর ভূমিকা ব্যাখ্যা করো।
\end{enumerate}\end{multicols}
